\chapter{Gestión}
\label{sec:gestion}

Este capítulo describe los elementos organizacionales que influyen en la calidad del software de SpeechNotes.

\section{Organización}
\label{sec:organizacion}

El equipo SQA está compuesto por 4 personas con roles definidos, manteniendo independencia funcional respecto al equipo de desarrollo. La estructura organizacional se muestra en la Figura~\ref{fig:organigrama}.

\begin{figure}[H]
\centering
\fbox{\parbox{0.8\textwidth}{\centering
\textbf{Project Manager}\\
$\downarrow$\\[-0.2cm]
\textbf{Consultora de Calidad (Externa)}\\
$\swarrow$ \quad $\downarrow$ \quad $\searrow$\\[-0.2cm]
\begin{tabular}{ccc}
\textbf{Pers. A} & \textbf{Pers. B} & \textbf{Pers. C}\\
QA Lead & QA Backend & QA Frontend\\
& + CI/CD & \\
\end{tabular}\\
$\downarrow$\\[-0.2cm]
\textbf{Pers. D}\\
QA Manual
}}
\caption{Organigrama del equipo SQA}
\label{fig:organigrama}
\end{figure}

\section{Responsabilidades}
\label{sec:responsabilidades}

La tabla~\ref{tab:responsabilidades} detalla las responsabilidades asignadas a cada miembro del equipo.

\begin{longtable}{cp{5cm}p{7cm}}
\caption{Matriz de responsabilidades SQA}
\label{tab:responsabilidades}\\
\toprule
\textbf{Rol} & \textbf{Responsable} & \textbf{Tareas Clave}\\
\midrule
\endfirsthead
\toprule
\textbf{Rol} & \textbf{Responsable} & \textbf{Tareas Clave}\\
\midrule
\endhead
\bottomrule
\endfoot

QA Lead & Persona A & Configuración de Jira, redacción del SQAP (Sección A), informe de cierre (Sección C), conclusiones (Sección D), ensamblaje del PDF final.\\

QA Backend & Persona B & Pipeline SonarCloud + GitHub Actions, pruebas automatizadas pytest, refactorización de code smells, reporte de análisis estático antes/después.\\

QA Frontend & Persona C & Configuración de Cypress/Jest, automatización de flujos clave (grabación, transcripción), corrección de bugs frontend reportados.\\

QA Manual & Persona D & Diseño de matriz de rastreabilidad, casos de prueba manuales (14 CP), ejecución manual, reporte de bugs en Jira, re-testing, video demo.\\

\end{longtable}

\section{Recursos}
\label{sec:recursos}

\subsection{Equipo y ambiente de prueba}
\label{subsec:ambiente}

\begin{itemize}
    \item \textbf{Sistema Operativo:} Windows 11.
    \item \textbf{Navegador:} Google Chrome / Chromium.
    \item \textbf{Frontend:} \url{http://localhost:3006}
    \item \textbf{Backend:} \url{http://localhost:9443}
    \item \textbf{Base de datos:} SQLite (dev.db) + MongoDB (Docker).
    \item \textbf{Repositorio:} GitHub, rama \texttt{planning-SQA}.
\end{itemize}

\subsection{Herramientas}
\label{subsec:herramientas}

\begin{itemize}
    \item \textbf{Gestión de incidencias:} Jira Software (Cloud).
    \item \textbf{Análisis estático:} SonarCloud.
    \item \textbf{CI/CD:} GitHub Actions.
    \item \textbf{Pruebas backend:} pytest.
    \item \textbf{Pruebas frontend:} Cypress + Jest.
    \item \textbf{Consola API:} curl / Postman.
    \item \textbf{Evidencias:} Capturas de pantalla, logs de backend, consola del navegador.
\end{itemize}
